\section{Model and Preliminaries}

We model each individual as being described by a tuple $((x, x'), y)$,
where $x\in \cX$ denotes a vector of \emph{protected attributes},
$x'\in \cX'$ denotes a vector of \emph{unprotected attributes}, and
$y\in \{0, 1\}$ denotes a label. Note that in our formulation, an
auditing algorithm not only may not see the unprotected attributes
$x'$, it may not even be aware of their existence. For example, $x'$
may represent proprietary features or consumer data purchased by a
credit scoring company.

We will write $X = (x, x')$ to denote the joint feature vector. We
assume that points $(X, y)$ are drawn i.i.d. from an unknown
distribution $\mathcal{P}$.  Let $D$ be a decision making algorithm,
and let $D(X)$ denote the (possibly randomized) decision induced by
$D$ on individual $(X,y)$. We restrict attention in this paper to the
case in which $D$ makes a binary classification decision:
$D(X) \in \{0,1\}$. Thus we alternately refer to $D$ as a classifier.
When \emph{auditing} a fixed classifier $D$, it will be helpful to
make reference to the distribution over examples $(X, y)$ together
with their induced classification $D(X)$. Let $\PA(D)$ denote the
induced \emph{target joint distribution} over the tuple $(x, y, D(X))$
% joint distribution} over the tuple $(x, x', y, D(X))$
that results from sampling $(x, x', y) \sim \mathcal{P}$, and
providing $x$, the true label $y$, and the classification
$D(X) = D(x,x')$ but not the unprotected attributes $x'$.  Note that
the randomness here is over both the randomness of $\mathcal{P}$, and
the potential randomness of the classifier $D$.


We will be concerned with learning and auditing classifiers $D$
satisfying two common statistical fairness constraints: equality of
classification rates (also known as statistical parity), and equality
of false positive rates (also known as equal opportunity). Auditing
for equality of false negative rates is symmetric and so we do not
explicitly consider it. Each fairness constraint is defined with
respect to a set of protected groups. We define sets of protected
groups via a family of indicator functions $\cF$ for those groups,
defined over protected attributes. Each $g:\cX \to \{0,1\} \in \cF$
has the semantics that $g(x) = 1$ indicates that an individual with
protected features $x$ is in group $g$.


\begin{definition}[Statistical Parity (SP) Subgroup Fairness]
  Fix any classifier $D$, distribution $\mathcal{P}$, collection of
  group indicators $\cG$, and parameter $\gamma\in [0, 1]$.
  For each $g\in \cG$, define
  \begin{align*}
    \spsize(g, \cP) = \Pr_{\cP}[g(x) = 1] \quad \mbox{and,} \quad
    \spdisp(g, D, \cP) = \left|\SP(D) - \SP(D, g)\right|,
  \end{align*}
  where $\SP(D) = \Pr_{\cP, D}[D(X) = 1]$ and
  $\SP(D, g) = \Pr_{\cP, D}[D(X) = 1 | g(x) = 1]$ denote the overall
  acceptance rate of $D$ and the acceptance rate of $D$ on group $g$
  respectively.  We say that $D$ satisfies $\gamma$-\emph{statistical
    parity (SP) Fairness} with respect to $\mathcal{P}$ and $\cF$ if
  for every $g \in \cF$
  \[ \spsize(g, \cP) \; \spdisp(g, D, \cP)\leq \gamma.\] We will
  sometimes refer to $\SP(D)$ as the \emph{SP base rate}.
\end{definition}

\begin{remark}
  Note that our definition references two approximation parameters,
  both of which are important. We are allowed to ignore a 
  group $g$ if it (or its complement) represent only a small fraction
  of the total probability mass. The parameter $\alpha$ governs how
  small a fraction of the population we are allowed to
  ignore. Similarly, we do not require that the probability of a
  positive classification in every subgroup is exactly equal to the
  base rate, but instead allow deviations up to $\beta$. Both of these
  approximation parameters are necessary from a statistical estimation
  perspective. We control both of them with a single parameter
  $\gamma$.
\end{remark}

\begin{definition}[False Positive (FP) Subgroup Fairness]\label{fp-fair}
  Fix any classifier $D$, distribution $\mathcal{P}$, collection of
  group indicators $\cF$, and parameter $\gamma \in [0,1]$.  For each
  $g\in \cG$, define
  \begin{align*}
    \fpsize(g, \cP) = \Pr_{\cP}[g(x) = 1, y=0] \quad \mbox{and,} \quad
                      \fpdisp(g, D, \cP) = \left| \FP(D) - \FP(D, g) \right|
  \end{align*}
  where $\FP(D) = \Pr_{D, \cP} [D(X) = 1\mid y=0]$ and
  $\FP(D, g) = \Pr_{D, \cP}[D(X) = 1\mid g(x) =1 , y=0]$ denote the
  overall false-positive rate of $D$ and the false-positive rate of
  $D$ on group $g$ respectively.


  We say $D$
  satisfies $\gamma$-\emph{False Positive (FP) Fairness} with respect
  to $\mathcal{P}$ and $\cF$ if for every $g \in \cG$
  \[ \fpsize(g, \cP) \; \fpdisp(g, D, \cP) \leq \gamma.\] We will
  sometimes refer to $\FP(D)$ FP-base rate.% , which we write as:
  % $b_{FP} = b_{FP}(D, P) = \Pr[D(X) = 1 | y = 0].$
\end{definition}

\begin{remark}
  This definition is symmetric to the definition of statistical parity
  fairness, except that the parameter $\alpha$ is now used to exclude
  any group $g$ such that \emph{negative examples} ($y = 0$) from
  $g$ (or its complement) have probability mass less than
  $\alpha$. This is again necessary from a statistical estimation perspective.
\end{remark}

For either statistical parity and false positive fairness, if the
algorithm $D$ fails to satisfy the $\gamma$-fairness condition, then
we say that $D$ is $\gamma$-\emph{unfair} with respect to
$\mathcal{P}$ and $\cF$. We call any subgroup $g$ which witnesses this
unfairness an $\gamma$-\emph{unfair certificate} for
$(D, \mathcal{P})$.

%\begin{definition}[Notions of Subgroup Fairness]\label{def:fair}
%  Fix any algorithm $D$, distribution $P$, and class of group
%  classifiers $\cF$. Consider the following three notions of
%  \emph{base rates}:
%\begin{align*}
%  b_1 &= b_1(D, P) = \Pr[D(X) = 1]\\
%  b_2 &= b_2(D, P) = \Pr[D(X) = 1 \mid y = 0]\\
%  b_3 &= b_3(D, P) = \Pr[y = 1 \mid D(X) = 1]
%\end{align*}
%For any $\alpha, \beta \in (0, 1)$, let
%$\cF_\alpha = \{g\in \cF \mid \Pr[g(x) = 1], \Pr[g(x) = 0] \geq
%\alpha\}$, and we say $D$ satisfies
%\begin{itemize}
%\item $(\alpha, \beta)$-\emph{statistical parity (SP) fairness} w.r.t. $P$ and $\cF$ if for
%  every $g\in \cF_\alpha$%  such that
%  % $\Pr[f(x) = 1], \Pr[f(x) = 0] \geq \alpha$
%  :
%  \[ | \Pr[D(X) = 1 | g(x) = 1] - b_1 | \leq \beta
%  \]
%\item $(\alpha, \beta)$-\emph{false-positive (FP) fairness} w.r.t. $P$ and
%  $\cF$ if for every $g\in
%  \cF_\alpha$: % such that $\Pr[f(x) = 1], \Pr[f(x) = 0] \geq \alpha$
%  \[ | \Pr[D(X) = 1 | g(x) = 1, y = 0] - b_2 | \leq \beta
%  \]
%\item $(\alpha, \beta)$-\emph{positive predictive (PP) fairness}
%  w.r.t. $P$ and $\cF$ if for every $g\in \cF_\alpha$:
%  \[
%    | \Pr[y = 1 | g(x) = 1, D(X) = 1] - b_3 | \leq \beta.
%  \]
%\end{itemize}
%For any fixed fairness notion above, if the algorithm $D$ fails to
%meet the $(\alpha, \beta)$-fairness condition, then we say $D$ is
%$(\alpha, \beta)$-\emph{unfair} w.r.t. $P$ and $\cF$, and we call the
%corresponding subgroup $f$ to be an $(\alpha, \beta)$-\emph{unfair
%  certificate} of $(D, P)$.
%\end{definition}




\iffalse
\begin{definition}[Auditable Class]
  Fix any of the fairness notions defined above. We say that a
  function class $F$ is \emph{(efficiently) auditable} if there exists
  an auditing algorithm $\cT$ such that for any
  $\alpha, \beta \in (0, 1)$, algorithm $D$ and any distribution $P$
  over $(x, x', y)$, when given access to labeled examples $(x, D(X))$
  from the underlying distribution, with probability $(1-\delta)$,
  $\cT$ does the following: If $D$ is $(\alpha, \beta)$-unfair
  w.r.t. $P$ and $F$, then $\cT$ will output the unfair certificate of
  $D$ under $\poly(1/\alpha, 1/\beta, 1/\delta)$ time;
    % such that
  % \[
  %   \Pr_P[f(x) = 1] \geq \alpha \qquad \mbox{ and, } \qquad |
  %   \Pr[D(x,x') = 1 | f(x) = 1] - b(D, P) | \leq \beta;
  % \]
  % \item If $D$ is $(\alpha/2, \beta/2)$-fair w.r.t. $P$ and $F$, then
  %   $\cT$ will output ``fair'' under time
  %   $\poly(1/\alpha, 1/\beta, 1/\delta)$ time.

\end{definition}
\fi

An \emph{auditing algorithm} for a notion of fairness is given sample access to $\PA(D)$ for some classifier $D$. It will either deem $D$ to be fair with respect to $\mathcal{P}$, or will else produce a certificate of unfairness.



\begin{definition}[Auditing Algorithm]
  Fix a notion of fairness (either statistical parity or false
  positive fairness), a collection of group indicators $\cF$ over the
  protected features, and any $\delta, \gamma, \gamma' \in (0, 1)$
  such that $\gamma'\leq \gamma$.  A {$(\gamma, \gamma')$-auditing
    algorithm for $\cF$ with respect to distribution $\mathcal{P}$} is
  an algorithm $\cA$ such that for any classifier $D$, when given
  access the distribution $\PA(D)$, $\cA$ runs in time
  $\poly(1/\gamma', \log(1/\delta))$, and with probability $1-\delta$,
  outputs a $\gamma'$-unfair certificate for $D$ whenever $D$ is
  $\gamma$-unfair with respect to $\mathcal{P}$ and $\cF$. If $D$ is
  $\gamma'$-fair, $\cA$ will output ``fair''.
  % $D$ is $(\alpha, \beta - 2\rho)$-fair w.r.t. $P$ and $\cF$,
  % then $A$ will output ``fair.''
\end{definition}

\iffalse
\begin{definition}[Auditing]
  Fix a notion of fairness (either statistical parity or
  false-positive fairness), a collection of group indicators $\cF$
  over the protected features, and any $\gamma, \gamma' \in (0, 1]$
  such that $\gamma'\leq \gamma$.  A collection of classifiers
  $\mathcal{H}$ is \emph{$(\gamma, \gamma')$-(efficiently) auditable
    under distribution $\mathcal{P}$ with respect to $\cF$} if there
  exists an auditing algorithm $A$ such that for every classifier
  $D \in \mathcal{H}$, when given access the distribution $\PA(D)$,
  $A$ runs in time $\poly(1/\gamma', \log(1/\delta))$, and with
  probability $(1-\delta)$, outputs an $\gamma'$-unfair certificate
  for $D$ whenever $D$ is $\gamma$-unfair with respect to
  $\mathcal{P}$ and $\cF$.
  % \item If $D$ is $(\alpha, \beta - 2\rho)$-fair w.r.t. $P$ and $\cF$,
  %   then $A$ will output ``fair.''
\end{definition}
\fi
As we will show, our definition of auditing is closely related to  weak
agnostic learning.

\begin{definition}[Weak Agnostic Learning~\citep{agnostic,KMV08}]
  Let $Q$ be a distribution over $\cX\times\{0, 1\}$ and let
  $\eps, \eps' \in (0, 1/2)$ such that $\eps \geq \eps'$. We say that
  the function class $\cF$ is \emph{$(\eps , \eps')$-weakly
    agnostically learnable} under distribution $Q$ if there exists an
  algorithm $L$ such that when given sample access to $Q$, $L$ runs in
  time $\poly(1/\eps', 1/\delta)$, and with probability $1 - \delta$,
  outputs a hypothesis $h\in \cF$ such that
  \[
    \min_{f\in\cF} err(f, Q) \leq 1/2 - \eps \implies err(h, Q) \leq
    1/2 - \eps'.
  \]
  % where $M = \min\{err_P(\mathbf{1}), err_P(\mathbf{0})\}$ denotes the
  % the smaller error achieved by the constant functions.
  where $err(h,Q) = \Pr_{(x,y) \sim Q}[h(x) \neq y]$.
\end{definition}
% \sw{make err function consistent throughout}

\paragraph{Cost-Sensitive Classification.}
In this paper, we will also give reductions to \emph{cost-sensitive classification
  (CSC)} problems. Formally, an instance of a CSC problem for the
class $\cH$ is given by a set of $n$ tuples
$\{(X_i, c_i^0, c_i^1)\}_{i=1}^n$ such that $c_i^\ell$ corresponds to
the cost for predicting label $\ell$ on point $X_i$. Given such an
instance as input, a CSC oracle finds a hypothesis $\hat h \in \cH$ that
minimizes the total cost across all points:
\begin{equation}
  \hat h \in \argmin_{h\in \cH} \sum_{i = 1}^n [h(X_i) c_i^1 + (1 -
  h(X_i)) c_i^0]\label{csc}
\end{equation}
A crucial property of a CSC problem is that the solution is invariant
to translations of the costs.



\begin{claim}
\label{claim:csc}
  Let $\{(X_i, c_i^0, c_i^1)\}_{i=1}^n$ be a CSC instance, and
  $\{(\tilde c_i^0, \tilde c_i^1)\}$ be a set of new costs such that
  there exist $a_1, a_2,\ldots, a_n\in \RR$ such that
  $\tilde c_i^\ell = c_i^\ell + a_i$ for all $i$ and $\ell$. Then
\[
  \argmin_{h\in \cH} \sum_{i = 1}^n [h(X_i) c_i^1 + (1 - h(X_i))
  c_i^0] = \argmin_{h\in \cH} \sum_{i = 1}^n [h(X_i) \tilde c_i^1 + (1
  - h(X_i)) \tilde c_i^0]
\]
\end{claim}

\begin{remark}
  We note that cost-sensitive classification is polynomially
  equivalent to agnostic learning~\cite{ZLA03}.  We give both
  definitions above because when describing our results for auditing,
  we wish to directly appeal to known hardness results for weak
  agnostic learning, but it is more convenient to describe our
  algorithms via oracles for cost-sensitive classification.
\end{remark} % \ar{Added this}

\paragraph{Follow the Perturbed Leader.}

We will make use of the \emph{Follow the Perturbed Leader (FTPL)}
algorithm as a no-regret learner for online linear optimization
problems \citep{KV05}.  To formalize the algorithm, consider
$\cS \subset \{0, 1\}^d$ to be a set of ``actions'' for a learner  in an online decision problem. The
learner interacts with an adversary over $T$ rounds, and in each round
$t$, the learner (randomly) chooses some action $a^t \in \cS$, and the
adversary chooses a loss vector $\ell^t \in [-M , M]^d$. The learner incurs a loss of $\langle \ell^t, a^t \rangle$ at round $t$.

FTPL is a simple algorithm that in each round perturbs the cumulative
loss vector over the previous rounds
$\overline \ell = \sum_{s< t} \ell^s$, and chooses the action that
minimizes loss with respect to the perturbed cumulative loss vector. We present the full
algorithm in \Cref{alg:ftpl}, and its formal guarantee in
\Cref{thm:ftpl-regret}.


\begin{algorithm}[h]
  \caption{Follow the Perturbed Leader (FTPL) Algorithm}
 \label{alg:ftpl}
  \begin{algorithmic}
    \STATE{\textbf{Input}: Loss bound $M$, action set
      $\cS \in \{0, 1\}^d$ } \STATE{\textbf{Initialize}: Let
      $\eta = (1/M) \sqrt{\frac{1}{\sqrt{d} T}} $, $\cD_U$ be the
      uniform distribution over $[0, 1]^d$, and let $a^1\in \cS$ be
      arbitrary.}

    \STATE{\textbf{For} $t = 1, \ldots, T$:}
    \INDSTATE{Play action
      $a^t$; Observe loss vector $\ell^t$ and suffer loss
      $\langle \ell^t , a^t \rangle$.}

    \INDSTATE{Update:}
    \[
      a^{t+1} = \argmin_{a\in \cS} \left[ \eta \sum_{s \leq t} \langle
        \ell^s , a \rangle + \langle \xi^t , a \rangle\right]
    \]
    \INDSTATE{where $\xi^t$ is drawn independently for each $t$ from the distribution $\cD_U$.  }
    \end{algorithmic}
  \end{algorithm}


  \begin{theorem}[\cite{KV05}]\label{thm:ftpl-regret}
    For any sequence of loss vectors $\ell^1, \ldots , \ell^T$, the
    FTPL algorithm has regret
    \[
      \Ex{}{\sum_{t=1}^T \langle \ell^t , a^t \rangle } - \min_{a\in
        \cS} \sum_{t=1}^T \langle \ell^t, a \rangle \leq 2 d^{5/4}
      M\sqrt{T}
      %% Calculations:
      %% L = 1 , sigma = sqrt{d} , D = sqrt{d}  , G = sqrt{d} M
      %  regret \leq 2 L D G sqrt{sigma T}
      %% eta^2 = sigma/ (G^2 L T) = 1 / (\sqrt{d} M^2 T)
      %% eta = sqrt{1 / \sqrt{d} T} * (1/M)
    \]
    where the randomness is taken over the perturbations $\xi^t$ across
    rounds.
  \end{theorem}

% \ar{There are probably theorems that get better dependencies on $d$ if we care.}




\subsection{Generalization Error}
\label{subsec:gen}
In this section, we observe that the error rate of a classifier $D$,
as well as the degree to which it violates $\gamma$-fairness (for both
statistical parity and false positive rates) can be accurately
approximated with the empirical estimates for these quantities on a
dataset (drawn i.i.d. from the underlying distribution $\mathcal{P}$)
so long as the dataset is sufficiently large. Once we establish this
fact, since our main interest is in the computational problem of
auditing and learning, in the rest of the paper, we assume that we
have direct access to the underlying distribution (or equivalently,
that the empirical data defines the distribution of interest), and do
not make further reference to sample complexity or overfitting issues.



A standard VC dimension bound (see, e.g. \cite{KV94}) states:
\begin{theorem}
  Fix a class of functions $\cH$. For any distribution $\mathcal{P}$,
  let $S \sim \mathcal{P}^m$ be a dataset consisting of $m$ examples
  $(X_i, y_i)$ sampled i.i.d. from $\mathcal{P}$. Then for any
  $0 < \delta < 1$, with probability $1-\delta$, for every
  $h \in \cH$, we have:
$$\left|err(h, \mathcal{P}) - err(h, S) \right| \leq O\left(\sqrt{\frac{\mathrm{VCDIM}(\cH)\log m + \log(1/\delta)}{m}} \right)$$
where
$err(h, S) = \frac{1}{m}\sum_{i=1}^m \mathbbm{1}[h(X_i) \neq y_i]$.
\end{theorem}

The above theorem implies that so long as $m \geq \tilde O(\mathrm{VCDIM}(\cH)/\epsilon^2)$, then minimizing error over the empirical sample $S$ suffices to minimize error up to an additive $\epsilon$ term on the true distribution $\mathcal{P}$. Below, we give two analogous statements for fairness constraints:
\begin{theorem}[SP Uniform Convergence]
\label{thm:SP-uc}
Fix a class of functions $\cH$ and a class of group indicators
$\cG$. For any distribution $\mathcal{P}$, let $S \sim \mathcal{P}^m$
be a dataset consisting of $m$ examples $(X_i, y_i)$ sampled
i.i.d. from $\mathcal{P}$. Then for any $0 < \delta < 1$, with
probability $1-\delta$, for every $h \in \cH$ and $g \in \cG$
\[
  \left| \spsize(g, \cP_S) \; \spdisp(g, h, \cP_S) - {\spsize(g, \cP) \;
      \spdisp(g, h, \cP)}\right| \leq
  \tilde{O}\left(\sqrt{\frac{(\mathrm{VCDIM}(\cH) +
        \mathrm{VCDIM}(\cG))\log m + \log(1/\delta)}{m}} \right)
\]
where $\cP_S$ denotes the empirical distribution over the realized
sample $S$.
\end{theorem}

Similarly:
\begin{theorem}[FP Uniform Convergence]
\label{thm:FP-uc}
Fix a class of functions $\cH$ and a class of group indicators
$\cG$. For any distribution $\mathcal{P}$, let $S \sim \mathcal{P}^m$
be a dataset consisting of $m$ examples $(X_i, y_i)$ sampled
i.i.d. from $\mathcal{P}$. Then for any $0 < \delta < 1$, with
probability $1-\delta$, for every $h \in \cH$ and $g \in \cG$, we
have:
$$
\left| \fpsize(g, \cP) \; \fpdisp(g, D, \cP) - {\fpsize(g, \cP) \;
    \fpdisp(g, D, \cP)}\right| \leq
\tilde{O}\left(\sqrt{\frac{(\mathrm{VCDIM}(\cH) +
      \mathrm{VCDIM}(\cG))\log m + \log(1/\delta)}{m}} \right)$$ where
$\cP_S$ denotes the empirical distribution over the realized sample
$S$.
\end{theorem}

These theorems together imply that for both SP and FP subgroup
fairness, the degree to which a group $g$ violates the constraint of
$\gamma$-fairness can be estimated up to error $\epsilon$, so long as
$m \geq \tilde O((\mathrm{VCDIM}(\cH) +
\mathrm{VCDIM}(\cG))/\epsilon^2)$. The proofs can be found in Appendix
\ref{app:generalization}.




%\paragraph{Remark}
%If the function class contains the constant 0 and constant 1
%functions, then the weak learner outputs meaningful hypothesis only if
%the optimal error $err_P(h^*)$ is no more than $(1/2 -\gamma)$.

% \paragraph{Problem}
% Suppose we are given points $(x, x')$ drawn from $P$ along with the
% algorithm's decisions $D(x, x')$. Is there any efficient algorithm (in
% terms of both sample complexity and computational complexity) that
% detects whether the algorithm satisfies $(\alpha, \beta)$-fairness or
% not
%%% Local Variables:
%%% mode: latex
%%% TeX-master: "main"
%%% End:
